\documentclass[12pt, a4paper]{article}

\usepackage{amsmath, amsthm, amssymb}
\usepackage[margin=2.5cm]{geometry}
\usepackage{hyperref}

% ── Theorem environments ──────────────────────────────────────────────────────
\theoremstyle{plain}
\newtheorem{theorem}{Theorem}
\newtheorem{lemma}{Lemma}

\theoremstyle{remark}
\newtheorem{remark}{Remark}

% ── Document ──────────────────────────────────────────────────────────────────
\begin{document}

\begin{center}
    {\LARGE \textbf{The Fundamental Theorem of Arithmetic}}\\[0.4em]
    {\large A proof via Euclid's Lemma}
\end{center}

\vspace{1em}

\begin{remark}
$1$ is not a prime number.
\end{remark}

\vspace{0.5em}

% ── Euclid's Lemma ───────────────────────────────────────────────────────────
\begin{lemma}[Euclid's Lemma]
If a prime $p$ divides the product $ab$ of two integers $a$ and $b$, then $p$
must divide at least one of those integers $a$ or $b$.
\end{lemma}

\vspace{0.5em}

% ── Fundamental Theorem of Arithmetic ────────────────────────────────────────
\begin{theorem}[Fundamental Theorem of Arithmetic]
Every natural number greater than $1$ is either a prime number or can be
expressed as a product of prime numbers. This representation is unique, except
for the order in which the factors appear.
\end{theorem}

\vspace{0.5em}

% ── Proof of Existence ────────────────────────────────────────────────────────
\begin{proof}[Proof of Existence]
We proceed by induction on $n$.

\medskip
\textbf{Base case} ($n = 2$).
Since $2$ is a prime number, the statement holds.

\medskip
\textbf{Induction hypothesis.}
Assume that every integer between $2$ and $n$ is either prime or a product of
primes.

\medskip
\textbf{Induction step.}
Let $m = n + 1$. We show that $m$ also satisfies the statement. There are two
cases.

\begin{enumerate}
  \item \emph{$m$ is a prime number.} The statement holds immediately.

  \item \emph{Otherwise,} there are integers $a$ and $b$ where $m = a \cdot b$
        and $1 < a \leq b < m$. By the induction hypothesis,
        $a = p_1 \cdot p_2 \cdots p_j$ and $b = q_1 \cdot q_2 \cdots q_k$
        are products of primes. But then
        \[
            m \;=\; a \cdot b
              \;=\; p_1 \cdot p_2 \cdots p_j \cdot q_1 \cdot q_2 \cdots q_k
        \]
        is also a product of primes, completing the induction.
        \qedhere
\end{enumerate}
\end{proof}

\vspace{0.5em}

% ── Proof of Uniqueness ───────────────────────────────────────────────────────
\begin{proof}[Proof of Uniqueness]
Suppose, for the sake of contradiction, that there exists an integer that has
two distinct prime factorizations. Let $n$ be the least such integer, and write
\[
    n \;=\; p_1 \cdot p_2 \cdots p_j \;=\; q_1 \cdot q_2 \cdots q_k,
\]
where each $p_i$ and each $q_i$ is prime.

\medskip
We see that $p_1$ divides $q_1 \cdot q_2 \cdots q_k$, so $p_1$ divides some
$q_i$ by Euclid's Lemma. Without loss of generality, say $p_1$ divides $q_1$.
Since $p_1$ and $q_1$ are both prime, it follows that $p_1 = q_1$.

\medskip
Returning to our factorizations of $n$, we may cancel these two factors to
conclude that
\[
    p_2 \cdots p_j \;=\; q_2 \cdots q_k.
\]
We now have two distinct prime factorizations of some integer strictly smaller
than $n$, which contradicts the minimality of $n$.
\end{proof}

\end{document}
